\subsection{The Poisson formula}

In this section, we consider a closed subgroup H of G. We denote \(\beta = dx\) the Haar measure on G and \(\widehat{\beta}\) the dual Haar measure on \(\widehat{G}\) . One denotes by \(\alpha\) a Haar measure on H and \(\widehat{\alpha}\) the dual Haar measure on the dual group \(\widehat{H}\) , which we identify with \(\widehat{G}/H^{\perp}\) (Theorem 4 of

II, p.~226). We also identify \(\widehat{\mathrm{G}}/\widehat{\mathrm{H}}\) with \(\mathrm{H}^{\perp}\) via the dual mapping of the canonical projection \(\mathrm{G}\to\mathrm{G}/\mathrm{H}\) (loc. cit.).

We shall denote by \(\dot{x}\) the canonical image of an element \(x\) of G in G/H and by \(\dot{\chi}\) the canonical image of an element \(\chi\) of \(\widehat{G}\) in \(\widehat{G}/H^{\perp}\) .

We denote by \(\gamma\) the Haar measure \(\beta/\alpha\) on G/H (INT, VII, §2, no. 2, def. 1 and no. 7, prop. 10), and \(\widehat{\gamma}\) the dual Haar measure on \(H^{\perp}\). Recall (INT, VII, §2, n° 3, prop. 5, c)) that the measure \(\gamma\) is characterized by the following property: for every \(f \in \mathcal{L}^{1}(G)\), the function \(y \mapsto f(xy)\) on H is \(\alpha\)-integrable for \(\beta\)-almost all \(x \in G\); its integral depends only on \(\dot{x}\) and the function defined \(\gamma\)-almost everywhere on G/H by

\[
f ^ {\flat} \colon \dot {x} \mapsto \int_ {\mathrm{H}} f (x h) d \alpha (h)
\]

belongs to L \(^{1}\) (G/H, γ) and satisfies

\[
\int_ {\mathrm{G} / \mathrm{H}} f ^ {\flat} d \gamma = \int_ {\mathrm{G}} f d \beta .\tag{33}
\]

PROPOSITION 15. --- Let \(f \in \mathrm{L}^1(\mathrm{G})\) such that the restriction to \(\mathrm{H}^\perp\) of the continuous function \(\mathcal{F}_{\mathrm{G}}(f)\) is integrable with respect to \(\widehat{\gamma}\). Then, for almost all \(x \in \mathrm{G}\), the function \(y \mapsto f(xy)\) on \(\mathrm{H}\) is \(\alpha\)-integrable, and one has:

\[
\int_ {\mathrm{H}} f (x y) d \alpha (y) = \int_ {\mathrm{H} ^ {\perp}} \langle \chi , x \rangle \mathcal {F} _ {\mathrm{G}} (f) (\chi) d \widehat {\gamma} (\chi).
\]

Based on the foregoing, the function \(f^{\flat}\) defined almost everywhere on G/H by

\[
f ^ {\flat} (\dot {x}) = \int_ {\mathrm{H}} f (x y) d \alpha (y)
\]

belongs to \(L^{1}(G/H)\) . The Fourier transform of \(f^{\flat}\) is identified with the function on \(H^{\perp} = \widehat{G/H}\) given for \(\chi \in H^{\perp}\) by

\[
\begin{array}{r} \mathcal {F} _ {\mathrm{G/H}} (f ^ {\flat}) (\chi) = \int_ {\mathrm{G/H}} \overline {{\langle \chi , \dot {x} \rangle}} f ^ {\flat} (\dot {x}) d \gamma (\dot {x}) \\ = \int_ {\mathrm{G}} \overline {{\langle \chi , x \rangle}} f (x) d \beta (x) = \mathcal {F} _ {\mathrm{G}} (f) (\chi) \end{array}
\]

by formula (33), applied to the integrable function \(x \mapsto \overline{\langle\chi, x\rangle} f(x)\) . By hypothesis, the function \(\mathcal{F}(f)|\mathrm{H}^{\perp} = \mathcal{F}_{\mathrm{G}/\mathrm{H}}(f^{\flat})\) belongs to \(\mathrm{L}^{1}(\mathrm{H}^{\perp})\) , and hence the function \(f^{\flat}\) belongs to the space B(G/H). It follows (th. 3 of II, p.~222) that \(f^{\flat}\) coincides almost everywhere with \(\overline{\mathcal{F}}_{\widehat{\mathrm{G / H}}}(\mathcal{F}_{\mathrm{G / H}}(f^{\flat}))\) . For almost every \(\dot{x}\in \mathrm{G / H}\) , we therefore have

\[
f ^ {\flat} (\dot {x}) = \int_ {\mathrm{H} ^ {\perp}} \langle \chi , x \rangle \mathcal {F} _ {\mathrm{G/H}} (f ^ {\flat}) (\chi) d \widehat {\gamma} (\chi) = \int_ {\mathrm{H} ^ {\perp}} \langle \chi , x \rangle \mathcal {F} _ {\mathrm{G}} (f) (\chi) d \widehat {\gamma} (\chi).
\]

This concludes the proof.

COROLLARY (Poisson formula). --- Let \(f \in \mathcal{L}^{1}(G)\) . It is assumed that the following conditions are satisfied:

\begin{enumerate}
\def\labelenumi{(\roman{enumi})}
\item
  The restriction of \(\mathcal{F}_{\mathrm{G}}(f)\) to \(\mathrm{H}^{\perp}\) is integrable;
\item
  For every \(x \in G\) , the function \(y \mapsto f(xy)\) on H is integrable;
\item
  The map \(x\mapsto \int_{\mathrm{H}}f(xy)d\alpha (y)\) is continuous on G. Then one has
\end{enumerate}

\[
\int_ {\mathrm{H}} f (y) d \alpha (y) = \int_ {\mathrm{H} ^ {\perp}} \mathcal {F} _ {\mathrm{G}} (f) (\chi) d \widehat {\gamma} (\chi).\tag{34}
\]

Indeed, taking up the notation of the proof of the preceding proposition, the functions \(f^{\flat}\) and \(\overline{\mathcal{F}}_{\widehat{\mathrm{G/H}}}\left(\mathcal{F}_{\mathrm{G/H}}(f^{\flat})\right)\) on G/H are continuous and equal almost everywhere. They are therefore equal everywhere and in particular at e, which gives formula (34).

PROPOSITION 16. --- The measure \(\widehat{\alpha}\) on \(\widehat{\mathrm{H}} = \widehat{\mathrm{G}} / \mathrm{H}^{\perp}\) is equal to \(\widehat{\beta} / \widehat{\gamma}\) .

Fix \(f\in \mathcal{K}(\mathrm{G})\) nonzero. For \(x\in \mathrm{G}\) and \(\chi \in \widehat{\mathrm{G}}\) , set

\[
\varphi (x, \chi) = \int_ {\mathrm{H}} f (x y) \langle \chi , y \rangle d \alpha (y).
\]

The function \(\varphi\) is continuous on \(G \times \widehat{G}\) (INT, IV, §4, no. 3, cor. 1 of th. 2). For \(x\) fixed, \(\varphi(x, \chi)\) depends only on the class of \(\chi\) in \(\widehat{G}/H^{\perp} = \widehat{H}\) . For \(\chi\) fixed, \(\langle \chi, x \rangle \varphi(x, \chi)\) depends only on the class of \(x\) in \(G/H\) , and the function \(\dot{x} \mapsto \langle \chi, x \rangle \varphi(x, \chi)\) on \(G/H\) has compact support.

Let \(x \in G\) . The function \(\dot{\chi} \mapsto \varphi(x, \chi)\) on \(\hat{H}\) is the Fourier cotransform of the function \(y \mapsto f(xy)\) on H. The latter is square-integrable, so by the Plancherel formula (23) of II, p.~217, we have

\[
\int_ {\widehat {\mathrm{G}} / \mathrm{H} ^ {\perp}} | \varphi (x, \chi) | ^ {2} d \widehat {\alpha} (\dot {\chi}) = \int_ {\mathrm{H}} | f (x y) | ^ {2} d \alpha (y).\tag{35}
\]

 Let \(\chi\in\widehat{\mathrm{G}}\) . The function \(\dot{x}\mapsto\langle\chi,x\rangle\varphi(x,\chi)\) belongs to \(\mathcal{H}(\mathrm{G}/\mathrm{H})\) , hence to \(\mathrm{L}^{1}(\mathrm{G}/\mathrm{H})\) . Its Fourier cotransform is the function on \(\mathrm{H}^{\perp}\)

whose value at \(\xi\in\mathrm{H}^{\perp}\) is

\[
\begin{array}{r l} & {\int_ {\mathrm{G/H}} \langle \xi , \dot {x} \rangle \langle \chi , x \rangle \varphi (x, \chi) d \gamma (\dot {x}) = \int_ {\mathrm{G/H}} \Bigl (\int_ {\mathrm{H}} \langle \chi \xi , x y \rangle f (x y) d \alpha (y) \Bigr) d \gamma (\dot {x})} \\ & {\qquad = \int_ {\mathrm{G}} \langle \chi \xi , x \rangle f (x) d \beta (x) = \overline {{\mathcal {F}}} _ {\mathrm{G}} (f) (\chi \xi)} \end{array}
\]

by formula (33). Therefore

\[
\int_ {\mathrm{G/H}} | \varphi (x, \chi) | ^ {2} d \gamma (\dot {x}) = \int_ {\mathrm{H} ^ {\perp}} | \overline {{\mathcal {F}}} _ {\mathrm{G}} (f) (\chi \xi) | ^ {2} d \widehat {\gamma} (\xi)\tag{36}
\]

by the Plancherel formula again.

One then finally computes

\[
\begin{array}{l l} \int_ {\widehat {G}} | \overline {{\mathcal {F}}} _ {G} (f) | ^ {2} d \widehat {\beta} = \int_ {G} | f | ^ {2} d \beta & (\text {by (23)}) \\ = \int_ {G / H} d \gamma (\dot {x}) \int_ {H} | f (x y) | ^ {2} d \alpha (y) & (\text {by (33)}) \\ = \int_ {G / H} d \gamma (\dot {x}) \int_ {\widehat {G} / H ^ {\perp}} | \varphi (x, \chi) | ^ {2} d \widehat {\alpha} (\dot {\chi}) & (\text {by (35)}) \\ = \int_ {\widehat {G} / H ^ {\perp}} d \widehat {\alpha} (\dot {\chi}) \int_ {G / H} | \varphi (x, \chi) | ^ {2} d \gamma (\dot {x}) \\ = \int_ {\widehat {G} / H ^ {\perp}} d \widehat {\alpha} (\dot {\chi}) \int_ {H ^ {\perp}} | \overline {{\mathcal {F}}} _ {G} (f) (\chi \xi) | ^ {2} d \widehat {\gamma} (\xi) & (\text {by (36)}) \end{array}
\]

where INT, V, §8, no. 3, prop. 5 has been applied to the continuous positive function \((\dot{x},\dot{\chi})\mapsto |\varphi (x,\chi)|^{2}\) on \(\mathrm{G / H}\times \widehat{\mathrm{G}} /\mathrm{H}^{\perp}\) .

Comparing this equality with the integration formula (33) for the group \(\widehat{G}\) , one concludes that the Haar measures \(\widehat{\alpha}\) and \(\widehat{\beta}/\widehat{\gamma}\) coincide.

